\section{Introduction}\label{sec:introduction}
\subsection{Motivation and scope}
Forecasts of electrical load provide an input to operational planning: a forecast error changes the amount of demand that must be accommodated beyond the anticipated load. This project studies the prediction problem itself; it does not simulate dispatch, electricity prices or reserve costs. Its practical relevance therefore rests on improving a clearly defined load forecast, rather than on claiming unmeasured economic benefits.

The Tetouan dataset provides an instructive setting because daily cycles, weekday differences, warm-season peaks and special-day patterns coexist within a single year. Strong dependence on yesterday's load also makes a simple seasonal naive forecast a demanding reference. A high coefficient of determination alone can conceal a model that adds little beyond that reference. Furthermore, records separated by ten minutes are strongly related: a random split of such observations can reward interpolation between neighboring times without answering the intended future-forecast question.

The implemented study uses the UCI Power Consumption of Tetouan City dataset\cite{uci} and forecasts the sum of its three zone loads after hourly averaging. The main setting, called \hday\ in the repository, assembles predictions for a day's hourly loads at midnight using daily lags. A secondary \hone\ setting uses more recent completed-hour observations. Both settings use the weather measured at the target hour, so their results are conditional on an oracle-weather assumption. This qualification is central to interpreting the experiment, not an incidental implementation detail.

\subsection{Research objectives and questions}
The repository's staged design supports three principal questions and two supplementary investigations:
\begin{enumerate}
\item How accurately do conventional ML regressors predict the hourly city total relative to seasonal naive forecasts under chronological evaluation?
\item How do calendar information, special-day indicators and historical load groups affect predictive performance?
\item How much do hyperparameter search and the target formulation change accuracy, temporal stability and recorded computational cost?
\item How do results change when recent hourly history becomes available in the secondary horizon setting?
\item How do shuffled evaluation and changes in the forecasting task alter apparent accuracy?
\end{enumerate}
These questions are answered using the saved experiments, not additional model development conducted while preparing this report.

\subsection{Contributions and evidence boundaries}
The contribution is an auditable comparison of five model families---linear regression, decision tree, random forest, support vector regression and XGBoost---within a shared forecasting pipeline. Its main components are common seasonal baselines; five expanding-window validation folds; randomized search over absolute and lag-relative targets; fixed-parameter feature ablations; a separately tuned hourly setting; and a diagnostic study of evaluation choices. The final analysis retains the model selected by development CV even when a different candidate has a lower observed December error.

Several boundaries prevent stronger claims. The dataset contains one year from one city. Feature decisions and parameter searches reuse development folds. Full-year descriptive graphics expose December's broad behavior, and the progress log records later reruns of the final notebook. The implementation establishes chronological fitting and a CV-dependent selection rule, but does not establish that the test data were hidden from every human decision throughout development. These distinctions are maintained throughout the report.

\subsection{Organization}
\Cref{sec:related,sec:data,sec:features,sec:method} establish the literature context, dataset, feature pipeline and evaluation protocol. \Cref{sec:baseline,sec:cv,sec:tuning,sec:ablation} examine validation and development experiments. \Cref{sec:final,sec:leakage} analyze the December results and diagnostic evaluation settings. \Cref{sec:discussion,sec:conclusion} synthesize the findings and identify future work. Supplementary material records exact source mappings, environment information and selected configurations.
\evidence{Guideline PDF pp. 1--5; literature-review PDF pp. 25--32; notebooks 02--08; \code{report/research_audit.md}.}
